% !TeX program = xelatex
% ============================================================
% vkr/vkr.tex — Выпускная квалификационная работа
% ОП «Программная инженерия» ФКН НИУ ВШЭ.
% Оформление: ГОСТ 7.32-2017 «Отчёт о научно-исследовательской
% работе. Структура и правила оформления». Класс extreport — главы.
%
% Структура (методичка ОП ПИ + ГОСТ 7.32):
%   Титульный лист → Реферат → Abstract → Содержание →
%   Термины и определения → Список сокращений → Введение →
%   3 главы → Заключение → Список источников → Приложения.
%
% Как пользоваться: пишите свой текст вместо пояснений (строки,
% начинающиеся с «%», в PDF не попадают). Данные о себе, научном
% руководителе, группе и годе задайте в настройках проекта —
% они подставятся в титул и реферат автоматически (макросы \hse…).
% ============================================================
\documentclass[12pt,a4paper]{extreport}
\input{preamble}
\begin{document}
\input{title}

% ============================================================
\chapter*{Реферат}
\addcontentsline{toc}{chapter}{Реферат}
% ============================================================
% Реферат (ГОСТ 7.32-2017 §5.3) — краткое и точное изложение работы.
% Структура одного абзаца (150–250 слов): проблема и актуальность →
% цель → что сделано (требования, проектирование, реализация) →
% главный результат и область его применения.

Здесь приводится краткое изложение работы: проблема и её
актуальность, цель работы, выполненные задачи
((* if project.kind == 'research' *))(аналитический обзор состояния проблемы, разработка
модели и метода решения, экспериментальная проверка результатов)((* else *))(анализ предметной
области, формирование требований, проектирование и реализация)((* endif *)),
а также главный полученный результат и область его применения.

\medskip
\textbf{Работа содержит} \pageref{LastPage}~страниц, \ldots~глав,
\ldots~рисунков, \ldots~таблиц, \ldots~источников и \ldots~приложения.

\medskip
% Ключевые слова (ГОСТ 7.32-2017 §5.3.2.1, §6.12.2): от 5 до 15 слов или
% словосочетаний, ПРОПИСНЫМИ буквами, в именительном падеже, в строку
% через запятую, без точки в конце.
\textbf{Ключевые слова:} ПРОГРАММНАЯ ИНЖЕНЕРИЯ, \ldots

\clearpage

% ============================================================
\chapter*{Abstract}
\addcontentsline{toc}{chapter}{Abstract}
% ============================================================
% English abstract (150–250 words) — обязателен для ВКР ФКН НИУ ВШЭ
% даже если работа написана на русском. Это перевод реферата.

This section contains a concise English summary of the thesis: the
problem and its relevance, the objective, the work carried out
((* if project.kind == 'research' *))(a systematic review of the state of the art, the proposed
model and method, and experimental evaluation)((* else *))(domain analysis, requirements,
design and implementation)((* endif *)), and the
main result with its area of application.

\medskip
\textbf{The work contains} \pageref{LastPage}~pages, \ldots~chapters,
\ldots~figures, \ldots~tables, \ldots~sources and \ldots~appendices.

\medskip
\textbf{Keywords:} SOFTWARE ENGINEERING, \ldots

\clearpage

% ── Содержание / Contents ───────────────────────────────────
\tableofcontents
\clearpage


% ============================================================
%  ENGLISH THESIS BODY (project.lang == en)
%  §2.5.3.1: an English VKR is written in English. The Реферат (Russian)
%  and Abstract (English) above are BOTH mandatory regardless of language.
%  References follow IEEE style (§2.5.3.2). Fill the yellow placeholders.
% ============================================================
\chapter*{Terms and Definitions}
\addcontentsline{toc}{chapter}{Terms and Definitions}
This thesis uses the following terms and their definitions.

\begin{description}[leftmargin=0pt, labelindent=0pt, font=\bfseries, style=nextline]
  \item[API] Application Programming Interface --- a set of conventions for
    interaction between software components.
  \item[Prototype] a partial implementation of the system built to validate
    key architectural and functional decisions.
  \item[\ldots] \ldots\ % add the remaining terms
\end{description}

\clearpage

\chapter*{List of Abbreviations and Symbols}
\addcontentsline{toc}{chapter}{List of Abbreviations and Symbols}
This thesis uses the following abbreviations and symbols.

\begin{description}[leftmargin=0pt, labelindent=0pt, font=\bfseries, style=nextline]
  \item[API] Application Programming Interface.
  \item[DBMS] Database Management System.
  \item[UI] User Interface.
  \item[\ldots] \ldots\ % add the remaining abbreviations
\end{description}

\clearpage

\chapter*{Introduction}
\addcontentsline{toc}{chapter}{Introduction}
% The introduction must cover: relevance, aim, objectives, object, subject,
% methods, novelty and practical value, and the structure of the work
% (recommended length — at least 500 words).

\section*{Relevance}
Justify why the problem matters now: what problem is addressed, whom it
affects, and why existing approaches are insufficient. Cite sources with
\verb|\cite|, e.g.~\cite{article}.

\section*{Aim and Objectives}
\textbf{Aim} --- \ldots\ % one sentence: what exactly is built

\textbf{Objectives:}
\begin{enumerate}[label=\arabic*), leftmargin=1.25cm]
((* if project.kind == 'research' *))
  \item carry out an analytical review of the literature and existing approaches;
  \item identify open problems and formulate the research task;
  \item develop a model (method, algorithm) that solves the stated task;
  \item run a computational experiment and analyse the results.
((* else *))
  \item analyse the domain and existing solutions;
  \item elicit functional and non-functional requirements;
  \item design the architecture and the data model;
  \item implement a prototype and test it.
((* endif *))
\end{enumerate}

\section*{Object and Subject}
\textbf{Object of study} --- \ldots

\textbf{Subject of study} --- \ldots

((* if project.kind == 'research' *))
\section*{Research Methods}
List the research methods (analytical review, formal modelling, statistical
analysis, computational experiment) and the tools used in the work.
((* else *))
\section*{Methods and Tools}
List the methods (literature analysis, comparison of analogues, modelling,
experiment) and the tools used in the work.
((* endif *))

\section*{Scientific Novelty and Practical Value}
Describe what is new in the work and where the result can be applied.

\section*{Structure of the Thesis}
The thesis consists of an introduction, the main body, a conclusion, the list
of references and appendices. Chapter~1 \ldots; Chapter~2 \ldots; Chapter~3
\ldots\ The work is formatted per GOST~7.32-2017~\cite{gost732}.

\clearpage

((* if project.kind == 'research' *))
% ============================================================
\chapter{LITERATURE REVIEW AND PROBLEM STATEMENT}
% ============================================================
\section{State of the Art}
Give an analytical review of the domain: key concepts, existing theoretical
results and approaches to the problem under study. Cite sources with
\verb|\cite|, e.g.~\cite{article}.

\section{Comparative Analysis of Existing Approaches}
Compare known methods (approaches, models) against criteria relevant to the
study; the comparison is given in Table~\ref{tab:approaches}. In the cells use
\verb|\cmpplus| (present), \verb|\cmpminus| (absent), \verb|\cmpplanned| (partial).

\begin{table}[H]
\caption{Comparison of existing approaches}
\label{tab:approaches}
\centering
\begin{tabular}{|>{\raggedright\arraybackslash}p{5cm}|c|c|c|}
\hline
\textbf{Criterion} & \textbf{Approach~A} & \textbf{Approach~B} & \textbf{Proposed} \\
\hline
Criterion~1 & \cmpplus  & \cmpminus  & \cmpplus \\
\hline
Criterion~2 & \cmpminus & \cmpplanned & \cmpplus \\
\hline
\end{tabular}
\end{table}

\section{Open Problems and Research Task}
Based on the review, identify the gaps in existing approaches, state the
research hypothesis and give a formal statement of the task (inputs, expected
result, constraints, quality criteria).

\section{Chapter Summary}
Summarise the review and justify the chosen research direction.

\clearpage

% ============================================================
\chapter{MODEL AND SOLUTION METHOD}
% ============================================================
\section{Formal Statement}
Introduce the notation and state the task rigorously: inputs, expected result,
constraints and quality criterion (e.g. an objective $f(x)\to\min$ or the
hypothesis being tested).

\section{Proposed Model (Method, Algorithm)}
Describe the proposed method (model, algorithm): its main steps, assumptions
and parameters. Explain formal derivations as you go.

\section{Justification of Correctness}
Provide a theoretical justification of the method: proofs of properties,
estimates of computational complexity or accuracy, or a substantive argument.

\section{Chapter Summary}
State the key theoretical results of the chapter.

\clearpage

% ============================================================
% The third chapter of a research thesis depends on the nature of the work
% (the arc of the SE academic thesis is NOT fixed):
%   • experimental study — keep "EXPERIMENTAL STUDY" (methodology, data,
%     metrics, results — as below);
%   • engineering/theoretical work without a computational experiment —
%     replace the chapter with "IMPLEMENTATION AND EVALUATION".
% ============================================================
\chapter{EXPERIMENTAL STUDY}
% ============================================================
\section{Experimental Methodology}
Describe the experimental plan: research questions, conditions and the
procedure used to test the hypothesis.

\section{Data and Metrics}
Describe the datasets used (size, source, preparation) and the metrics for
evaluating the quality of the results.

\section{Results and Analysis}
Present the results obtained (Table~\ref{tab:results}), compare them with the
baselines and assess the significance of the differences.

\begin{table}[H]
\caption{Experimental results}
\label{tab:results}
\centering
\begin{tabular}{|>{\raggedright\arraybackslash}p{5cm}|c|c|}
\hline
\textbf{Method} & \textbf{Metric~1} & \textbf{Metric~2} \\
\hline
Baseline method  & \ldots & \ldots \\
\hline
Proposed method  & \ldots & \ldots \\
\hline
\end{tabular}
\end{table}

\section{Discussion}
Interpret the results, state the limitations of the study and the conditions
under which the conclusions apply.

\section{Chapter Summary}
Summarise the results of the experimental study.

\clearpage

((* else *))
% ============================================================
\chapter{DOMAIN ANALYSIS}
% ============================================================
\section{Domain Analysis}
Describe the domain, its key concepts and the problem the work solves.

\section{Review of Existing Solutions}
Compare existing analogues against relevant criteria. Every table must be
referenced in the text (e.g. the comparison is given in
Table~\ref{tab:analogues}). In the cells use \verb|\cmpplus| (yes),
\verb|\cmpminus| (no), \verb|\cmpplanned| (partial/planned).

\begin{table}[H]
\caption{Comparison of existing solutions}
\label{tab:analogues}
\centering
\begin{tabular}{|>{\raggedright\arraybackslash}p{5cm}|c|c|c|}
\hline
\textbf{Criterion} & \textbf{Analogue~A} & \textbf{Analogue~B} & \textbf{Proposed} \\
\hline
Criterion~1 & \cmpplus  & \cmpminus  & \cmpplus \\
\hline
Criterion~2 & \cmpminus & \cmpplanned & \cmpplus \\
\hline
\end{tabular}
\end{table}

\section{System Requirements}
State the functional and non-functional requirements. If the requirements are
defined in the Statement of Work via \verb|\req{ТЗ-Ф-01}{...}|, reference them
in the text with \verb|\reqref{ТЗ-Ф-01}| --- the Requirements traceability
panel links them automatically. The functional requirements are listed in
Table~\ref{tab:func-req}.

\begin{table}[H]
\caption{Functional requirements}
\label{tab:func-req}
\centering
\begin{tabular}{|c|>{\raggedright\arraybackslash}p{11cm}|}
\hline
\textbf{ID} & \textbf{Description} \\
\hline
ТЗ-Ф-01 & \ldots \\
\hline
\end{tabular}
\end{table}

\section{Chapter Summary}
Summarise the analysis and justify the problem statement.

\clearpage

% ============================================================
\chapter{DESIGN}
% ============================================================
\section{Solution Architecture}
Describe the overall architecture. Every figure must be mentioned in the text
% Пример ссылки показываем ТЕКСТОМ: живой \ref{fig:arch} ссылался на метку,
% которой в шаблоне нет, — в PDF получалось «??», а нормоконтроль справедливо
% ругался на висячую ссылку. В русской редакции это уже сделано через \verb,
% но здесь пример стоит внутри аргумента \enquote{}, где \verb ломается.
(e.g. \enquote{the overall diagram is shown in Figure~\texttt{\textbackslash ref\{fig:arch\}}}). Insert
an image with \verb|\includegraphics| (a missing file is replaced by a
placeholder) or draw it with \texttt{tikz}.

\section{Data Model}
Describe the main entities and relationships (an ER diagram or a text
description).

\section{Technology Stack}
List the languages, frameworks and libraries with a justification of the choice.

\section{Chapter Summary}
State the design decisions made.

\clearpage

% ============================================================
\chapter{IMPLEMENTATION AND TESTING}
% ============================================================
\section{Implementation Details}
Describe the key modules and non-trivial decisions. Format code fragments with
the \texttt{lstlisting} environment (see Listing~\ref{lst:example}).

\begin{lstlisting}[language=Python, caption={Example function}, label={lst:example}]
def greet(name: str) -> str:
    return f"Hello, {name}!"
\end{lstlisting}

\section{Testing}
Describe the testing strategy (unit, integration, e2e) and give the results in
Table~\ref{tab:test-results}.

\begin{table}[H]
\caption{Testing results}
\label{tab:test-results}
\centering
\begin{tabular}{|>{\raggedright\arraybackslash}p{4cm}|>{\raggedright\arraybackslash}p{8cm}|c|}
\hline
\textbf{Test type} & \textbf{What is verified} & \textbf{Result} \\
\hline
Unit & \ldots & \ldots \\
\hline
Integration & \ldots & \ldots \\
\hline
\end{tabular}
\end{table}

\section{Evaluation and Results}
Describe where and how the solution was validated and the results obtained.

\section{Chapter Summary}
Summarise the implementation and testing.

\clearpage
((* endif *))

\chapter*{Conclusion}
\addcontentsline{toc}{chapter}{Conclusion}
% Brief results per objective from the introduction, the main result, its
% correspondence to the aim, and directions for further work.
As a result of the work:
\begin{enumerate}[label=\arabic*), leftmargin=1.25cm]
  \item \ldots\ % result for objective 1
  \item \ldots\ % result for objective 2
\end{enumerate}

The aim of the work has been achieved. Directions for further work: \ldots

\clearpage

% ── References (IEEE style, §2.5.3.2) ────────────────────────
% For an English thesis the reference list and in-text citations follow IEEE
% style: numbered [n], author initials first, article titles in quotes,
% venue/volume/number/pages, accessed date for online sources.
\phantomsection
\addcontentsline{toc}{chapter}{References}
\renewcommand{\refname}{References}
\begin{thebibliography}{99}

  \bibitem{book}
    A.~B.~Author, \emph{Title of the Book}, 2nd~ed. Moscow, Russia: Publisher,
    2023.

  \bibitem{article}
    A.~B.~Author and C.~D.~Coauthor, \enquote{Title of the article,}
    \emph{Journal Name}, vol.~12, no.~1, pp.~10--20, 2024.

  \bibitem{webresource}
    Resource Name. (2026). Accessed: Jan.~1, 2026. [Online]. Available:
    \url{https://example.org}

  \bibitem{gost732}
    \emph{GOST 7.32-2017. System of Standards on Information, Librarianship and
    Publishing. Research Report. Structure and Presentation Rules}. Moscow,
    Russia: Standartinform, 2017.

\end{thebibliography}

\clearpage

% ── Appendices (Latin letters A, B, C… for the English text) ──
\appendix
\renewcommand{\thechapter}{\Alph{chapter}}
\titleformat{\chapter}[hang]
  {\normalfont\large\bfseries\centering}{Appendix~\thechapter}{1ex}{}

\chapter{Source Code Listings}
\label{app:listing}
Place large listings, diagrams, data tables and other supporting material here.



\end{document}
